\section[The Part leaf]{The \code{Part} leaf}\label{sec:leaf}

A \code{Part} (\srcf{core/model/parts/Part.h}) answers exactly four questions about one physical
component: who it is (a stable id, a human-facing name, a type tag), what it weighs and how that
mass is distributed (own mass, a per-unit-mass inertia tensor, a CM offset), what silhouette it
presents along the axis (length and closed-form radius profiles), and what it contributes to the
Barrowman stack (a normalized \code{AeroComponent}). It answers nothing else. There is no parent
pointer, no child list, no placement, no cache, and no composite quantity anywhere on the class ---
the tree structure and everything derived from it belong to \code{PartNode} and \code{PartsModel}
(\cref{sec:tree}), and the derived state is cached there behind the two gates of
\cref{sec:caching}. The leaf is the passive, physical vocabulary; the tree is the grammar.

This division is enforced at the type level rather than by convention. A part is a
unique-ownership value: assignment is deleted, the copy constructor is protected, and duplication
exists only as the virtual \code{clone()}. Post-attach mutation exists only as two private setters
whose sole friend is \code{PartsModel}, so an edit that skips cache invalidation is not merely
discouraged --- it does not compile. \Cref{fig:leaf-class} shows the full contract and the five
concrete leaves; the rest of this section walks it top to bottom.

\classfig{class-part-leaf}{The \code{Part} leaf contract (\srcfcap{core/model/parts/Part.h}) and its
five concrete types. Concrete boxes list only the members each type overrides;
\code{PartsModel} appears only in its leaf-facing role, as the friend behind the routed
\code{setMass}/\code{setI} seam (and, for \code{Motor}, the motor verbs).}{fig:leaf-class}{1.0}

\subsection{Identity: id, name, type tag}\label{sec:leaf-identity}

Every part carries a \code{Part::Id} (\cppinline|std::uint64_t|), drawn from a process-wide atomic
counter at construction (\src{core/model/parts/Part.cpp}{17}). The counter starts at 1, reserving 0 as
the ``none/invalid'' value that other layers use for a root's parent or a cleared event field; it
increments with \cppinline|std::memory_order_relaxed| because uniqueness, not ordering, is the
contract. Every constructor assigns a fresh id --- including the protected copy constructor
(\src{core/model/parts/Part.cpp}{35}) --- so a clone is a new part, never an alias of its original.
Ids are deliberately not serialized: they are meaningful only within a single run, and a design
file round-trip (\cref{sec:serialization}) yields a structurally identical tree with new ids.

The \code{name} is a human-facing label with no uniqueness requirement;
\code{PartsModel::findByName} is a pre-order convenience, but identity is always the id. The type
tag \code{typeName()} is the one pure identity virtual: a stable string
(\code{"NoseCone"}, \code{"BodyTube"}, \code{"FinSet"}, \code{"HollowSphere"}, \code{"Motor"})
that doubles as the part-factory key (\srcf{core/model/parts/PartFactory.h}), the design-file
\code{<part type=...>} attribute, and the CLI \code{listparts} label. One vocabulary serves the
factory, the file format, and the front-ends, so the three cannot drift apart.

\subsection{Mass properties: the per-unit-mass split}\label{sec:leaf-mass}

The leaf stores its own mass in kilograms and its inertia tensor \emph{per unit mass} --- a purely
geometric tensor in $\unit{m^2}$, about the part's own CM (\src{core/model/parts/Part.h}{60}). The full
mass-weighted tensors in $\unit{kg\,m^2}$ exist only as composites, assembled by the node walk of
\cref{sec:tree}, which weights each part's geometric tensor by its live mass and shifts it to the
composite CM with the parallel-axis theorem.

\begin{listing}[htbp]
\begin{minted}{cpp}
/// Per-unit-mass (geometric) inertia tensor about this part's CM (m^2).
virtual Matrix3 getI() const { return inertiaTensor; }

/// CM relative to the component middle. Zero for symmetric parts, non-zero for a
/// cone (CM at L/4 or L/3 from the base). Consumed by the composite walk as
/// `-L/2 + getCenterMassOffset().z()`.
Vector3 getCenterMassOffset() const { return cm; }

/// This part's own mass at simulation time @p t (kg). Lets overrides model
/// time-varying mass (e.g. a motor).
virtual double getMass(double t [[maybe_unused]]) const
{
    return mass;
}
\end{minted}
\caption{The mass-property surface of the leaf (\srccap{core/model/parts/Part.h}{59}--70): a geometric
tensor, a middle-relative CM offset, and a time-parameterized mass.}\label{lst:leaf-mass}
\end{listing}

The split (\cref{lst:leaf-mass}) exists because mass is the only time-varying physical input in the model. The hook is
\code{getMass(t)}: the base implementation ignores $t$ and returns the stored mass, and
\code{Motor} overrides it to read its wrapped \code{MotorModel} live --- loaded weight before
ignition, falling to the casing mass over the burn (\cref{sec:motor}). Since the composite walk
re-weights resolved geometry by \code{getMass(t)} on every rebuild, this one virtual is what makes
every composite quantity --- mass, CM, inertia --- a function of time.

\begin{rejected}[Rejected --- mass-weighted leaf tensors]
Storing $\unit{kg\,m^2}$ at the leaf would bind a snapshot of the mass into the tensor. Every mass
edit would then require a paired tensor rescale --- two writes maintaining one invariant --- and a
burning motor would need its stored tensor rewritten every timestep. Factoring the mass out
entirely makes the two orthogonal: mass moves, geometry does not, and the tensor written at
construction stays valid for the part's whole life. The price is one scalar multiply per part per
composite rebuild, paid inside a walk that is already touching every part.
\end{rejected}

\begin{keyinsight}[Units at the leaf boundary]
\code{Part::getI()} is per-unit-mass, $\unit{m^2}$, about the part's own CM. Composite tensors
(\code{PartNode::compositeI}) are mass-weighted, $\unit{kg\,m^2}$, about the composite CM. The
weighting is applied exactly once --- \cppinline|c.mass * c.part->getI()| inside the node walk ---
so a leaf tensor is never double-weighted and a composite is never un-weighted.
\end{keyinsight}

The CM offset is deliberately expressed relative to the \emph{component middle}, not the fore
plane. The composite walk locates each part's CM at
\cppinline|-L/2 + getCenterMassOffset().z()| below its resolved fore-plane origin
(\src{core/model/PartsModel.cpp}{199}); middle-relative storage makes the zero vector the correct
value for every symmetric part (tube, sphere), so only genuinely asymmetric shapes record an
offset, computed once in their constructors. A fore-plane-relative convention would instead force
every symmetric part to store $-L/2$ and keep it synchronized with its length --- exactly the
derived-from-geometry data this design refuses to store. The offset is non-virtual and set once at
construction; its radial components are carried but unconsumed until 6-DOF force application.
Placement itself never reads the CM: attachment is geometric (\code{StationLink},
\cref{sec:placement}), and the CM participates only in the mass walk.

\subsection{The aerodynamic contribution}\label{sec:leaf-aero}

\code{getAero(refArea)} (\src{core/model/parts/Part.h}{75}) returns the part's Barrowman contribution
as an \code{AeroComponent} --- $C_{N_\alpha}$, the weighted moment \code{cnAlphaXcp}, and a drag
term --- normalized to the caller-supplied shared reference area so that contributions from
different parts are directly additive. The default is aerodynamically inert
(\cppinline|return {};|, $C_{N_\alpha} = 0$), and the value is a pure function of the part's
geometry: computed on demand, never stored.

\begin{designrule}[The own-CM aero datum]
Each part reports its center of pressure as $x_{\mathrm{cp}}$ measured from \emph{its own CM}
along $+z$. A leaf cannot know its station in the rocket --- that is placement, which lives in the
tree --- but its own CM is the one landmark both sides can locate: the part derives it from its
geometry, and the composite walk knows exactly where that CM lands
(\src{core/model/PartsModel.cpp}{237}). The walk adds $C_{N_\alpha}\cdot z_{\mathrm{cm}}$ to each
contribution, re-expressing every part on the sub-tree-root datum shared with
\code{compositeCm} --- which is what makes $\mathrm{cp} - \mathrm{cg}$ the static margin with no
frame conversion.
\end{designrule}

The component stores the product \code{cnAlphaXcp} rather than a raw $x_{\mathrm{cp}}$, so
composing a CP is plain field addition and a zero-lift body drops out of the weighted average with
no special case --- a \code{BodyTube} contributes $C_{N_\alpha} = 0$ and simply vanishes from the
sum (\srcf{core/model/Aero.h}).

\code{getReferenceArea()} is the second, decoupled half of the aero story: each part reports its
own frontal disc ($\pi r^2$, or 0 for a part with no frontal disc), and the tree's
\code{maxFrontalReferenceArea} takes the maximum over the sub-tree --- the shared Barrowman
reference area that is then fed back into every \code{getAero} call. A part thus \emph{reports} an
area once and \emph{receives} the winning area as a parameter; no leaf ever assumes it defines the
rocket's reference.

\subsection{The geometry profile}\label{sec:leaf-geometry}

Every part describes itself in the same local frame: it occupies $z \in [-L, 0]$ with its fore
plane at the origin and \zfwd{} (\src{core/model/parts/Part.h}{79}). \code{getLength()} defaults to 0,
which is the contract for a geometrically inert node --- a \code{Motor} contributes mass but
consumes no stack length --- and every concrete geometry type overrides it. The profile proper is
two closed-form virtuals, \code{radiusOuterAt(zLocal)} and \code{radiusInnerAt(zLocal)}: a tube
returns constants, a cone a linear taper, a sphere a circular arc. They are callable independently
of any station bookkeeping because the envelope sweep of \cref{sec:placement} samples them across
whole spans, not just at seams. \code{axialLength()} defaults to \code{getLength()} and exists as
the override point for a part whose envelope span differs from its nominal length; no current part
needs it, but the seam keeps ``envelope'' from silently aliasing ``length'' in the sweep.

\begin{listing}[htbp]
\begin{minted}{cpp}
Station Part::stationAt(double station01) const
{
    // Clamp the fraction (degenerate guard) and map it into the +z = forward local
    // frame: station 1 is the fore plane (z = 0, the origin), station 0 the aft
    // plane (z = -length).
    const double s = std::clamp(station01, 0.0, 1.0);
    const double z = (s - 1.0) * getLength();
    return Station{z, radiusOuterAt(z), radiusInnerAt(z)};
}

double Part::innerCapacityAt(double zLocal) const
{
    // Solid-host rule: a solid part is bounded by its outer skin (the poke-through
    // test); a bored part by its inner wall. Branching here keeps the overlap sweep
    // free of solidity special cases.
    return isSolid() ? radiusOuterAt(zLocal) : radiusInnerAt(zLocal);
}
\end{minted}
\caption{Fraction-to-landmark mapping and the solid-host rule
(\srccap{core/model/parts/Part.cpp}{44}--58).}\label{lst:leaf-station}
\end{listing}

\code{stationAt(station01)} (\cref{lst:leaf-station}) turns a dimensionless fraction --- 0 the aft
plane, 1 the fore plane --- into a resolved \code{Station} landmark: the axial coordinate plus the
outer and inner radii read from the profile there. The fraction is clamped to $[0,1]$, the base
mapping is $z = (s-1)\,L$, and because it is built entirely from \code{getLength()} and the radius
virtuals, no symmetric part needs an override; a \code{Station} is derived on every read and never
stored, so editing a length can never strand a stale landmark.

\code{isSolid()} is inferred, not stored: the base asks whether the part has a bore at its fore
plane (\cppinline|radiusInnerAt(0.0) <= 0.0|). Two classes of part must override the inference.
A \code{HollowSphere}'s cavity vanishes at the poles, so the fore-plane probe would misreport a
shell as solid; it answers from its inner radius directly. A shell \code{ConicalNoseCone} does not
model an inner profile at all, so its solidity is \emph{authored} --- the constructor's
\code{solid} flag is the answer. Getting this right matters because \code{isSolid()} selects the
edge in \code{innerCapacityAt}: the radius a host offers an inserted child is its bore if it has
one, and its outer skin (the poke-through bound) if it does not. Centralizing that branch in the
base class keeps the envelope sweep entirely free of per-type solidity cases --- an offender of
outer radius $r$ fits at a station iff $r \le$ \code{innerCapacityAt} $+$ tolerance, whatever the
host is.

\subsection[Duplication: clone() through a protected copy]{Duplication: \code{clone()} through a protected copy}\label{sec:leaf-clone}

A part is duplicated one way: \code{clone()}, a pure virtual returning
\cppinline|std::unique_ptr<Part>| (\src{core/model/parts/Part.h}{125}). Each concrete type implements
it as \cppinline|new BodyTube(*this)| against its defaulted copy constructor, which chains to
\code{Part}'s \emph{protected} copy constructor --- same mass properties, fresh id. The access
choreography is the point: because the base copy constructor is protected and both assignment
operators are deleted, external code cannot copy-construct a part at all, so slicing a
\code{BodyTube} down to a base \code{Part} is unrepresentable rather than a lurking bug. The
concrete types keep their own copy constructors and \code{clone()} overrides protected too, so
even a caller holding the concrete type goes through the virtual. \code{PartNode::clone} builds on
this to deep-copy whole sub-trees --- every leaf cloned with fresh ids, links copied verbatim
(\cref{sec:tree}).

\subsection{The routed-mutation seam}\label{sec:leaf-seam}

After a part is attached to the tree, exactly two things about it may change: its mass and its
tensor. Both writes are private, and their sole friend is \code{PartsModel}
(\cref{lst:leaf-seam}). The public paths are the routed verbs
\code{PartsModel::setPartMass} and \code{setPartInertia} (\src{core/model/PartsModel.cpp}{399}),
which perform the plain write and then dirty the mass-cache chain up the tree --- write and
invalidation fused in the only code that can do either.

\begin{listing}[htbp]
\begin{minted}{cpp}
protected:
    /// Copy for concrete clone() implementations only: same mass properties,
    /// fresh id. Protected so external code can't copy or slice.
    Part(const Part&);

private:
    friend class model::PartsModel;  // the sole post-attach mutation seam

    /// Plain writes: PartsModel's routed verbs own cache invalidation.
    void setMass(double m) { mass = m; }
    void setI(const Matrix3& I) { inertiaTensor = I; }
\end{minted}
\caption{The copy and mutation seams of the leaf (\srccap{core/model/parts/Part.h}{127}--137).}
\label{lst:leaf-seam}
\end{listing}

\begin{rejected}[Rejected --- protected setters]
Protected access would admit every subclass to the seam, and a subclass can re-publish what it
inherits: one public \cppinline|void setDensity(double)| wrapper that calls a protected
\code{setMass} internally, and the invalidation contract is silently gone --- the caches of
\cref{sec:caching} would keep serving mass properties for a part that no longer has them. Private
$+$ \code{friend} names the one class that may write, and that class is the one that owns cache
invalidation, so the seam cannot be widened by inheritance. \code{Motor} repeats the pattern one
level down: its \code{setMotorModel} is itself private with \code{PartsModel} as friend, so even a
motor swap is a routed verb (\cref{sec:motor}).
\end{rejected}

\subsection{The concrete leaves}\label{sec:leaf-concrete}

Five concrete types cover the current vocabulary; \cref{tab:leaf-overrides} maps each against the
virtuals it overrides. The umbrella header \srcf{core/model/parts/Parts.h} registers them for use
sites, and \code{makePart}/\code{params} (\srcf{core/model/parts/PartFactory.h}) construct and reflect
the geometric four from \code{typeName()} plus named parameters (\cref{sec:serialization,sec:cli}).

\begin{table}[htbp]
\centering\small
\begin{tabular}{@{}lccccc@{}}
\toprule
\textbf{Virtual} & \textbf{BodyTube} & \textbf{NoseCone} & \textbf{FinSet} & \textbf{Sphere} & \textbf{Motor}\\
\midrule
\code{typeName()} (pure)     & $\bullet$ & $\bullet$ & $\bullet$ & $\bullet$ & $\bullet$\\
\code{clone()} (pure)        & $\bullet$ & $\bullet$ & $\bullet$ & $\bullet$ & $\bullet$\\
\code{getLength()}           & $\bullet$ & $\bullet$ & $\bullet$ & $\bullet$ & ---\\
\code{radiusOuterAt()}       & $\bullet$ & $\bullet$ & $\bullet$ & $\bullet$ & ---\\
\code{radiusInnerAt()}       & $\bullet$ & --- & --- & $\bullet$ & ---\\
\code{isSolid()}             & $\bullet$ & $\bullet$ & --- & $\bullet$ & ---\\
\code{getReferenceArea()}    & $\bullet$ & $\bullet$ & $\bullet$ & --- & ---\\
\code{getAero()}             & $\bullet$ & $\bullet$ & $\bullet$ & --- & ---\\
\code{getMass(t)}            & --- & --- & --- & --- & $\bullet$\\
\code{getI()}                & --- & --- & --- & --- & $\bullet$\\
\code{axialLength()}         & --- & --- & --- & --- & ---\\
\code{stationAt()}           & --- & --- & --- & --- & ---\\
\bottomrule
\end{tabular}
\caption{Virtual overrides by concrete part. No type overrides \code{axialLength()} or
\code{stationAt()}: every current envelope equals its nominal length, and the base
fraction-to-landmark mapping serves all profiles. \code{Motor} overrides no geometry at all --- it
is the mass-only, zero-length leaf.}\label{tab:leaf-overrides}
\end{table}

\paragraph{BodyTube} (\srcf{core/model/parts/BodyTube.h}) is a uniform-density hollow cylinder:
$V = \pi(r_o^2 - r_i^2)L$, tensor from \code{InertiaTensors::Tube}, CM at the geometric center
(zero offset). Its profile is two constants --- outer skin and bore --- and $r_i = 0$ degrades it
to a solid rod, which its \code{isSolid()} reports directly. Aerodynamically it is the canonical
zero-lift body: a constant-diameter cylinder produces no normal force, so \code{getAero} returns
the inert component and the tube drops out of the composite CP average; its wetted area is
exposed for an eventual skin-friction drag term.

\paragraph{ConicalNoseCone} (\srcf{core/model/parts/ConicalNoseCone.h}) is a straight cone, uniform
solid by default or a thin open-base lateral shell. It is the working example of the CM-offset
machinery: the centroid sits $L/4$ (solid) or $L/3$ (shell) forward of the base, so the
constructor bakes $\bar h - L/2$ --- that is, $-L/4$ or $-L/6$ from the middle --- into the base
class via \code{coneCmOffset}, on top of any user-supplied offset. Its \code{getAero} carries the
matching datum conversion (\cref{lst:leaf-cone}): the Barrowman cone CP sits $\tfrac{2}{3}L$ aft
of the tip regardless of solidity, which lands at $L/3 - \bar h$ from the CM --- $+L/12$ for a
solid cone, exactly 0 for a shell, whose CM and CP coincide. The outer profile is a linear taper
with a guard so a degenerate zero-length cone reads as a flat disc rather than evaluating $0/0$;
solidity is authored by the constructor flag, as \cref{sec:leaf-geometry} requires.

\begin{listing}[htbp]
\begin{minted}{cpp}
// CM relative to the component middle, in the +z = forward frame (tip at +L/2,
// base at -L/2). The centroid is hbar forward of the base, i.e. hbar - L/2 from
// the middle: -L/4 (solid), -L/6 (shell).
Vector3 ConicalNoseCone::coneCmOffset(double L, bool solid)
{
    const double hbar = solid ? (L / 4.0) : (L / 3.0);
    return Vector3{0.0, 0.0, hbar - L / 2.0};
}

model::AeroComponent ConicalNoseCone::getAero(double refArea) const
{
    // Barrowman cone: CNalpha = 2 at the base area (pi R^2), rescaled to refArea.
    // The cone CP is 2/3 L aft of the tip (shape-independent). Report x_cp from the
    // cone's own CM (the composite datum) in the +z = forward frame:
    //   x_cp(from CM) = -L/6 - (hbar - L/2) = L/3 - hbar  ( = +L/12 solid, 0 shell ).
    const double cnAlpha = 2.0 * (std::numbers::pi * baseRadius * baseRadius) / refArea;
    const double hbar = solid ? (length / 4.0) : (length / 3.0);
    const double xcpFromCm = length / 3.0 - hbar;
    return {cnAlpha, cnAlpha * xcpFromCm, 0.0};
}
\end{minted}
\caption{The cone's centroid offset and its own-CM aero datum conversion
(\srccap{core/model/parts/ConicalNoseCone.cpp}{49}--68, lightly elided).}\label{lst:leaf-cone}
\end{listing}

\paragraph{FinSet} (\srcf{core/model/parts/FinSet.h}) models $N \ge 3$ identical trapezoidal
flat-plate fins as \emph{one} analytic leaf rather than $N$ child parts. The tree shares a single
body-frame orientation and never rotates child tensors, so an azimuthally-arrayed component cannot
be represented as children at different clock angles; instead the $N$-fin rotate-and-sum is baked
into the set's own per-unit-mass tensor (\code{InertiaTensors::TrapezoidalFinSet}), which is
transversely isotropic for $N \ge 3$. $N < 3$ is truly anisotropic and unsupported --- the
constructor warns and applies the $N \ge 3$ approximation rather than throwing. Its axial extent
is the root chord (\cppinline|getLength() == rootChord|), and its silhouette is deliberately the
mounting disc: \code{radiusOuterAt} returns \code{bodyRadius}, with the fin-tip radius
$r_b + s$ exposed separately as extent (\code{getMaxRadius}). The same reasoning fixes its
reference area: \code{getReferenceArea()} reports the body disc $\pi r_b^2$, \emph{not} a
span-inflated disc, because that value competes in the tree's max-frontal-disc search --- fins are
lifting surfaces, not frontal area, and letting a fin tip win the search would corrupt the
normalization of every coefficient in the rocket. \code{getAero} is the full Barrowman fin-set
term --- interference factor $1 + r_b/(s + r_b)$, referenced to the body cross-section and
rescaled to the shared area --- with the CP converted from the root leading edge to the set's own
mass centroid, and a degenerate guard that reports an inert component for $r_b = 0$ (fins on the
axis make the $(s/d)^2$ term singular, and with no body there is no interference to model).

\paragraph{HollowSphere} (\srcf{core/model/parts/HollowSphere.h}) is a thick-walled shell:
$V = \tfrac{4}{3}\pi(r_o^3 - r_i^3)$, per-unit-mass tensor
$\tfrac{2}{5}(r_o^5 - r_i^5)/(r_o^3 - r_i^3)\,\mathbf{1}$, axial extent the full diameter
(pole to pole), centered at $z = -r_o$. It is the one part with genuinely curved closed forms ---
outer silhouette $\sqrt{r_o^2 - (z + r_o)^2}$ and an inner cavity profile that vanishes outside
the band $|z + r_o| \le r_i$ --- and it is the motivating case for the \code{isSolid()} override:
its bore is zero at the fore plane, so only the direct $r_i > 0$ answer classifies a shell
correctly.

\paragraph{Motor} (\srcf{core/model/parts/Motor.h}) is surveyed here and specified in
\cref{sec:motor}. It wraps a \code{MotorModel} by value and overrides exactly the two mass
virtuals: \code{getMass(t)} reads the wrapped motor live, and \code{getI()} derives a solid-cylinder
per-unit-mass tensor from the catalog diameter and length on every call (degrading to a
zero-tensor point mass when the catalog omits geometry). It overrides no geometry virtual at all
--- length 0, no silhouette, no aero --- so it participates in the composite as pure time-varying
mass at its placed station, and its runtime mutations (in-place swap, ignition) go through
\code{PartsModel}'s motor verbs via its own private-plus-friend seam.

Taken together, the leaves hold every fact that is local to one component and nothing that is not.
What a part weighs \emph{now}, where those masses sit relative to each other, and whether the
assembly self-intersects are questions about the tree --- and the tree is where this document goes
next (\cref{sec:tree}).
